% ============================================================================
% PROPOSED — RQ1 results section (STAGING DRAFT; not part of the build)
% ============================================================================
% Prepared 2026-09-26. Insert target in main.tex: §5.2, currently
%   \section{Telemetry Delivery Semantics (RQ1)}\label{sec:results_rq1_delivery}
% whose TODO comment this draft replaces (heading text and label stay as they
% are).
%
% SYNC NOTE 2026-09-28: main.tex §5.2 updated — the delivery paragraph was
% trimmed to reference the new tab:rq1_observation (delivery realisation and
% overload observability); this draft mirrors the paragraph, the new table,
% and the British-English alignments (artefacts, cancelled). Exhibit order in
% both files: observation table → arms table → edges table → figure.
%
% SYNC NOTE 2026-09-29: the observation table and the §5.1 outcomes bullet
% were trimmed to the two relevant metrics (delivered fraction, information
% age at decision); the detection-delay and missed-window columns were cut
% from reporting (the campaign record retains them; appendix graphs
% overload_detection_delay / missed_overload still exist). This draft
% mirrors the trim.
%
% SYNC NOTE 2026-09-29 (narrative pass): the results narrative received an
% interpretation pass in both files — aim-first framing, per-result
% interpretation, reframed expectation deviations, answer-shaped closing;
% all numbers unchanged.
%
% SYNC NOTE 2026-09-30: margin-effect passage rewritten for explicitness
% (exposure window defined; absorption margin and per-arm outcomes stated;
% causal closing) — mirror of main.tex.
%
% SYNC NOTE 2026-10-01: main.tex §5.2 — a pre-registered expectations row
% was added to tab:rq1_design (the design table lives only in main.tex); the
% four result exhibits were re-ordered to sit after the paragraphs that
% invoke them (ISCTE near-invocation rule). This draft mirrors the new
% exhibit order.
% Follow-up (same date): the replicated-ordering sentence gained an 'on
% usable capacity' gloss (mirrored here); the expectations row fix in
% tab:rq1_design ('permitted' → 'expected') is main-only, as the design
% table lives only in main.tex.
%
% SYNC NOTE 2026-10-05: §5.2 numbers re-verified against the cloud-vm run
% artifacts (reaction_timeline / info_age / phase_service_quality /
% per_node_stats, plus the campaign dataset and stats): D-arm episode-p95
% range corrected (30.7 → 31.4; raw minimum 31.43 s, sp_2); margin-sentence
% usable-capacity range aligned to the per-run mean-of-domains basis
% (52.6--63.6 → 55.1--61.5 s; the wider figure was the pooled per-domain
% range, dominated by the flagged delayed_3 run); robustness sentence now
% notes the two single-window recovery blips (C 4.2 s, D 5.8 s).
%
% NUMBER PROVENANCE (all values re-verified 2026-09-26 against the campaign
% record):
%   - tese/research_questions/rq1/rq1_conclusions.md (F1--F10, verdicts)
%   - docs/operation/testing/experiment/v3/rq1/analysis/rq1_campaign_summary.md
%     (per-run table, pre-registered edges, non-surge table)
%   - docs/operation/testing/experiment/v3/rq1/post_run_analysis.md §5.1
%     (margin effect: control-loop timings, processing-time spread)
%   - observability values (delivered fraction, information age, detection
%     delay, window coverage) recomputed from the 28 run bundles'
%     analysis/rq1_delivery/*.csv on cloud-vm, 2026-09-26.
%
% FIGURE:
%   Source: docs/operation/testing/experiment/v3/rq1/graphs/thesis/
%           episode_latency_p95.png
%   Copy to: tese/images/rq1_episode_latency_p95.png
%   (pk: copy first, then build. Do NOT use endpoint_latency_p95.png — that
%   variant plots the run-wide pooled p95, which dilutes the episode and does
%   not match the pre-registered metric.)
%
% AFTER INSERTING (verification protocol):
%   - grep old TODO form -> zero matches; grep "usable capacity was 28.5" ->
%     present; reload main.tex.
%   - Update tese/Notes/purpose_evidence_map.md and
%     tese/miscelineous/plan_claim_aligment_with_measured_evidence.md (B4, C3
%     close on insert; C4 = archive the campaign dataset locally — still an
%     open decision, see chat note 2026-09-26).
%
% OPEN DECISIONS (do not resolve in the text yet):
%   1. Data-chart figure pipeline: the thesis figure instructions currently
%      say "every figure is drawn in draw.io"; result charts have no drawio
%      genre. Recommendation: codify a carve-out for source-generated result
%      charts (analysis pipeline) once approved.
%   2. Optional appendix material kept out of the body: per-bucket no-masking
%      check, run-isolation/scale-down integrity closure, per-run matrices,
%      alternate graphs (completeness_vs_infoage, missed_overload,
%      overload_detection_delay, scale_reaction_latency, per-phase latency).
% ============================================================================


\section{Telemetry Delivery Semantics (RQ1)}\label{sec:results_rq1_delivery}

The campaign compared the three delivery semantics named in
Section~\ref{sec:research_questions} under the overload episode described in
Section~\ref{sec:experimental_setup}, with the latest-state class realised in
two lossy forms. Arm~A (event-preserving) hands every completed telemetry
window to the decision engine exactly once and in source order; arm~B
(delayed event-preserving) hands the same ordered windows with a fixed,
pre-registered delay of 30~s and no burst replay; arm~C (latest-state
polling) delivers only the most recent completed window at each 30~s poll;
and arm~D (sampled push) delivers every third completed window immediately
and drops the intermediate windows. The aggregation window, scaling policy,
routing policy, and workload were held constant across arms. Each arm was run
seven times, once in each of seven seed blocks shared by the four arms,
giving 28 runs in total, with each run treated as an experimental unit.

The results are presented in the order of the research question: what the
controller observed about the overload, how it responded by scaling, and the
quality of service users experienced during the demand shift.

The delivered evidence matched this design
(Table~\ref{tab:rq1_observation}): arms A and B received every completed
telemetry window, while the lossy arms received about one third of them; the
information age at decision varied with the delivery treatment. In effect,
the same overload reached the controller through four different views: A
observed every window as it closed; B observed the same stream one fixed
delay behind; C obtained only the window current at each poll; and D
received one window in three as it was published.

\begin{table}[htbp]
    \centering
    \caption{Delivery realisation and overload observability (mean over the
    seven replicates; delivered fraction as the observed range): fraction of
    telemetry windows delivered and mean information age at decision.}
    \label{tab:rq1_observation}
    \footnotesize
    \setlength{\tabcolsep}{3pt}
    \begin{tabular}{l c c}
        \toprule
        Arm & Delivered & Information age (s) \\
        \midrule
        A -- event-preserving & 1.000 & 3.5 \\
        B -- delayed (+30~s) & 1.000 & 33.3 \\
        C -- latest-state (poll 30~s) & 0.325--0.336 & 16.2 \\
        D -- sampled push (/3) & 0.328--0.336 & 11.7 \\
        \bottomrule
    \end{tabular}
\end{table}

The scaling response followed the order in which the evidence reached the
controller. Every arm detected the overload and scaled; what differed is when
the response became usable capacity. The median time from the onset of the
episode to the first new capacity serving traffic was 28.5~s in A, 59.6~s in
B, 79.6~s in C and 83.2~s in D (Table~\ref{tab:rq1_arms}). Capacity became
usable first under A, and about 31~s later under B --- close to the 30~s lag
the treatment applies to every delivered window. Under C and D it took nearly
three times as long. Because the scaling policy, routing policy, and workload
were identical across arms, the differences in this column are attributable
to the delivery treatment. The first two steps of the pre-registered ordering
were fully replicated on usable capacity: A < B and B < C in all seven blocks. C and D were not
cleanly separated, with D later at the median by 3.6~s (C < D in five of
seven blocks). Reaction claims use usable-capacity latency; the time to the
first scale-up decision is reported as descriptive only, because delivery
timing confounds it by construction.

\begin{table}[htbp]
    \centering
    \caption{Delivery-semantics arms: usable capacity, episode latency, and
    outcome classes (medians with min--max ranges over the seven replicates;
    delivered fraction and timeouts as observed). Usable capacity is the time
    from episode onset to first new capacity serving traffic; the episode p95
    is the served-basis endpoint latency p95 (per-run mean of the two
    domains).}
    \label{tab:rq1_arms}
    \footnotesize
    \setlength{\tabcolsep}{2pt}
    \begin{tabular}{l c c c c c}
        \toprule
        Arm & Delivered & Usable capacity (s) & Episode p95 (s) & Timeouts & Cancelled (\%) \\
        \midrule
        A -- event-preserving & 1.000 & 28.5 (25.1--32.9) & 3.7 (2.8--5.2) & 0 & 5.6 \\
        B -- delayed (+30~s) & 1.000 & 59.6 (55.1--61.5) & 10.8 (5.0--65.3) & 0 & 6.5 \\
        C -- latest-state (poll 30~s) & 0.325--0.336 & 79.6 (73.7--81.8) & 35.0 (27.0--39.6) & 0 & 5.4 \\
        D -- sampled push (/3) & 0.328--0.336 & 83.2 (75.1--88.1) & 34.5 (31.4--37.6) & 0 & 6.1 \\
        \bottomrule
    \end{tabular}
\end{table}

The user-visible cost tracked the reaction, and here the two failure modes
separated (Table~\ref{tab:rq1_edges}, Figure~\ref{fig:rq1_episode_p95}).
Loss of intermediate evidence was the heavier and more robust cost: with
about two thirds of the windows absent, the episode endpoint p95 reached
35.0~s in C and 34.5~s in D against 3.7~s in A, roughly a ninefold penalty,
and the contrast separated completely --- every A replicate lay below every
C and D replicate (Cliff's $\delta = -1.000$, exact $p = 0.0006$). Delivery
delay alone also degraded the outcome in aggregate, to 10.8~s against A's
3.7~s ($\delta = -0.959$, exact $p = 0.0012$), but not uniformly: two of the
seven delayed replicates collapsed to lossy-arm levels (65.3~s and 64.7~s),
while the other five remained close to the reference (5.0--12.6~s).

\begin{table}[htbp]
    \centering
    \caption{Pre-registered primary contrasts on the episode endpoint latency
    p95 (served basis; per-run mean of the two domains). Exact Mann--Whitney
    $p$-values (complete enumeration, $n = 7$ per arm); Cliff's $\delta$
    computed as $P(x > y) - P(x < y)$ with 95~\% bootstrap CI; direction =
    blocks (of seven) consistent with the predicted order. No multiplicity
    correction (four pre-registered edges).}
    \label{tab:rq1_edges}
    \begin{tabular}{l c c c c}
        \toprule
        Contrast & Medians (s) & Exact $p$ & Cliff's $\delta$ [95~\% CI] & Direction \\
        \midrule
        Delay A--B & 3.73 vs 10.77 & 0.0012 & $-0.959$ [$-$1.000, $-$0.755] & 7/7 \\
        Delay D--C & 34.50 vs 35.04 & 0.456 & $-0.265$ [$-$0.837, 0.388] & 2/7 \\
        Loss A--D & 3.73 vs 34.50 & 0.0006 & $-1.000$ [$-$1.000, $-$1.000] & 7/7 \\
        Loss B--C & 10.77 vs 35.04 & 0.209 & $-0.429$ [$-$1.000, 0.429] & 5/7 \\
        \bottomrule
    \end{tabular}

    {\footnotesize\noindent\textit{Note.} Direction counts the blocks
    consistent with each contrast's predicted order: A $<$ B and D $>$ C
    (delay edges), A $<$ D and B $<$ C (loss edges). The D $>$ C prediction
    reflects the flag, before the campaign, that the sampled-push arm could
    be strictly worst; its low count (2/7) is the deviation reported in the
    text.\par}
\end{table}

\begin{figure}[H]
    \centering
    \includegraphics[width=0.75\textwidth]{images/rq1_episode_latency_p95.png}
    \caption{Episode endpoint latency p95 by delivery semantics (served
    basis; per-run mean of the two domains; seven replicates per arm). Bars
    show the median across replicates, whiskers the min--max range, and dots
    the individual runs; the two highest delayed points are the replicates
    where the base fleet's absorption margin was consumed.}
    \label{fig:rq1_episode_p95}
\end{figure}

A post-campaign analysis of the run artefacts attributes the spread of the
delayed arm to the margin each semantics leaves to the base fleet rather than
to the control loop. Across all seven delayed runs the control loop behaved
identically: the decision evidence was 30~s old, the first scale-up decision
fell between 52.0~s and 60.4~s, and usable capacity between 55.1~s and
61.5~s. The user-visible outcome depends on the exposure window: the time the base
fleet must absorb the surge alone before the new capacity serves traffic.
That window is about 29~s in A, 55--62~s in B, and 78--83~s in C and D.

Whether that window is survived also depends on the testbed's run-to-run
variance, which is present in every arm and is not attributable to the
delivery treatment: the episode median compute processing time varied
roughly sixfold (16--94~ms) across runs of identical configuration. The
base fleet has only a limited absorption margin against the surge, and a
slow platform state exhausts it sooner than a fast one. A's exposure window
is short enough that the margin covers it in every run, so the variance
leaves no visible trace in its outcome. The window of C and D is long
enough that the margin is exhausted in every run, so both arms degrade in
all seven replicates. B sits at the boundary: the margin covers its window
in five of the seven replicates, while in the two that sample a slow state
the margin runs out and the arm collapses to lossy-arm levels. The delivery
semantics set the exposure window, and the platform state decides whether
the base fleet's absorption margin covers it.

Two pre-registered expectations were not reproduced, and both are reported
as such. The full ordering A < B < C $\approx$ D did not hold, because B
does not sit reliably between A and C: its collapsed replicates straddle C
(B--C $p = 0.209$, direction 5/7). The sampled-push arm was also not strictly
worst: C and D were statistically indistinguishable ($p = 0.456$), with the
point estimate marginally against C. Read together with the margin effect
described above, both deviations are informative rather than negative. The
collapsed delayed replicates show that the cost of delay is conditional on
the platform state that the arm samples, and the equivalence of C and D
shows that, once the intermediate windows are gone, whether the loss arises
from polling or from sampling leaves the cost unchanged. The expectation of
a single severity ordering across the two failure modes therefore did not
hold in full; what the campaign measured instead is two failure modes that
differ in kind: loss of intermediate evidence acts robustly, while delivery
delay acts conditionally on the margin it consumes. One delayed replicate
also showed a domain asymmetry above the campaign's threefold screening line
(domain~2 p95 21.1~s against 4.0~s in domain~1); it was the only run above
that line.

The cost is confined to the overload episode. Outside it the arms were
indistinguishable, with baseline, recovery and demand-drop p95 at about
1.05--1.08~s in every arm (bar two single-window recovery blips, one in~C at
4.2~s and one in~D at 5.8~s), and no run recorded a single timeout inside the
episode. The only timeouts in the campaign were the client-cap tail timeouts
during the demand drop (0.07--0.12~\% of offered requests); cancellation and
dropped shares were flat across arms at roughly 5--7~\%, with the small
completed-without-status class remaining largest in B (2.22~\% against
0.03--0.42~\% elsewhere). A per-bucket check (30~s buckets) confirmed that
the episode ordering is not an artefact of aggregating the episode into a
single percentile. On the evaluated testbed, the observation interface
therefore shapes the transient response to a demand shift, while leaving
steady-state quality untouched.

In sum, the results answer the question in the order it was posed: the
delivery semantics changed what the controller observed about the overload;
they changed when the response became usable capacity, with the lossy arms
taking nearly three times as long as the reference; and they changed the
quality of service users experienced, with an approximately ninefold penalty
in the episode endpoint p95 where intermediate evidence was lost, and a
delay cost that scales with the margin the semantics leave to the base
fleet. Under this workload, the two failure modes of the observation
interface were therefore distinguishable and not interchangeable: loss of
intermediate demand evidence was the robust driver of user-visible cost,
while delivery delay was costly in proportion to the margin it consumes.
